% from ega0-20-fr.tex
\oldpage[0\textsubscript{IV-1}]{147}
complexes de $C$-modules « à homotopie près » (c'est-à-dire où les morphismes
sont les classes d'homomorphismes homotopes) sont des invariants remarquables,
plus fins que le couple formé de $\Omega^1_{C/B}$ et de
$\Upsilon_{C/B/A}$. Lorsque $k$ est un corps premier, et $C$ une $k$-algèbre
formellement lisse (par exemple un anneau régulier de type fini sur une
extension de $k$ (cf.\ ($\mathbf{IV}$, 6.8.6))), on peut montrer que le
complexe $\mathrm{K}_\bullet(C/A/k)$ peut se décrire uniquement en faisant
intervenir $C$ et $A$ (à l'exclusion de $k$) : on exprime $C$ comme quotient
d'une algèbre de polynômes $B$ sur $A$ par un idéal $\mathfrak{Q}$, et on
considère le complexe $\mathrm{F}_\bullet(C/A)$ à 2 termes non nuls
\[
  \cdots\to0\to\mathfrak{Q}/\mathfrak{Q}^2\to
  \Omega^1_{B/A}\otimes_B C\to0\to\cdots
\]
(dont l'homologie coïncide bien avec celle de
$\mathrm{K}_\bullet(C/A/k)$ en vertu de \sref{0.20.6.23.1-fr}). Ces
complexes $\mathrm{F}_\bullet(C/A)$, qui du point de vue de l'algèbre
homologique jouent le rôle d'un fibré conormal pour $\operatorname{Spec}(C)$
au-dessus de $\operatorname{Spec}(A)$, tiendront une place importante dans
les chapitres de cet ouvrage consacrés à la dualité des faisceaux cohérents et
au théorème de Riemann--Roch.
\end{remark}

\subsection{Généralisations aux anneaux topologiques}
\label{subsection:0.20.7-fr}

\begin{env}[20.7.1]
\label{0.20.7.1-fr}
Il résulte aussitôt des définitions que si, dans \sref{0.20.5.2-fr} et
\sref{0.20.5.3-fr}, les anneaux $A$, $B$, $C$ sont supposés être des anneaux
topologiques et les homomorphismes d'anneaux $u$, $v$ continus, alors les
homomorphismes $u_{C/B/A}$ et $v_{C/B/A}$ sont \emph{continus} (il faut
naturellement prendre sur $\Omega^1_{B/A}\otimes_B C$ la topologie produit
tensoriel).

En outre, $u_{C/B/A}:\Omega^1_{C/A}\to\Omega^1_{C/B}$ est un
\emph{morphisme strict surjectif} de $C$-modules topologiques ; en effet, il en
est ainsi de l'homomorphisme canonique $C\otimes_A C\to C\otimes_B C$, compte
tenu de la définition de la topologie produit tensoriel ; par passage aux
quotients on en déduit que l'homomorphisme correspondant
$\mathrm{P}^1_{C/A}\to\mathrm{P}^1_{C/B}$ est un morphisme strict surjectif,
et $u_{C/B/A}$ est la restriction de ce dernier à $\Omega^1_{C/A}$.

Dans \sref{0.20.5.5-fr}, si $A'$ et $B$ sont des $A$-algèbres topologiques,
et si $B'$ et $\Omega^1_{B/A}\otimes_B B'$ sont munis des topologies produits
tensoriels, l'homomorphisme canonique \sref{0.20.5.5.1-fr} est un isomorphisme
topologique, en tenant compte de ce que $\mathrm{P}^1_{B/A}$ est somme directe
topologique de $B$ et de $\Omega^1_{B/A}$ \sref{0.20.4.8-fr}.
\end{env}

\begin{proposition}[20.7.2]
\label{0.20.7.2-fr}
Soient $u:A\to B$, $v:B\to C$ deux homomorphismes continus d'anneaux
topologiques. Pour que l'homomorphisme continu
$v_{C/B/A}:\Omega^1_{B/A}\otimes_B C\to\Omega^1_{C/A}$ soit formellement
inversible à gauche (cf.\ \sref{0.19.1.5-fr}), il faut et il suffit que $C$
soit une $B$-algèbre formellement lisse relativement à $A$
\sref{0.19.9.1-fr} (et \emph{a fortiori} il suffit que $C$ soit une
$B$-algèbre formellement lisse).
\end{proposition}

\begin{proof}
Dire que $v_{C/B/A}$ est formellement inversible à gauche signifie en effet,
vu que les topologies de $\Omega^1_{C/A}$ et $\Omega^1_{B/A}\otimes_B C$
sont moins fines que celles déduites de la topologie de $C$
\sref{0.20.4.5-fr}, que pour tout $C$-module \emph{discret} $L$, annulé par un
idéal ouvert de $C$, l'homomorphisme canonique
\[
  \operatorname{Hom.\ cont}_C(\Omega^1_{C/A},L)\to
  \operatorname{Hom.\ cont}_C(\Omega^1_{B/A}\otimes_B C,L)
\]
